% !TeX root = prueba01_version02.tex
% !TeX program = pdflatex
\documentclass[10pt,a5paper]{article}
\newif\ifmostrarSoluciones
\mostrarSolucionesfalse
\usepackage{estilo_apuntes}
\dosPorCarillatrue % Cambie a \dosPorCarillatrue para imprimir dos páginas por carilla.
\encabezadoprueba{Prueba 1}{23 de octubre de 2026}

\begin{document}
\begin{center}
  {\color{Azul}\large\bfseries Prueba 1}
\end{center}
\noindent\textbf{Indicaciones}
\begin{itemize}
  \item La duración de la prueba es de 1 hora y 30 minutos.
  \item Escriba el desarrollo de cada ejercicio de manera clara y ordenada.
  Justifique cada paso e indique las propiedades, restricciones o resultados que
  utilice. Una respuesta sin justificación recibirá únicamente una nota parcial.
  \item No se permite el uso de ningún dispositivo electrónico durante la
  evaluación.
  \item Cualquier intento de copia penalizará la evaluación con una calificación
  total de cero.
\end{itemize}

\medskip
\noindent\textbf{Ejercicios}
\begin{enumerate}
  \item \textbf{(3 puntos)} Para $a\in\mathbb{R}$, considere la función
  \[
    F_a\colon\mathbb{R}\to\mathbb{R},
    \qquad
    F_a(x)=
    \begin{cases}
      \dfrac{x^2+ax-2a-4}{|x-2|}, & x\neq2,\\[2mm]
      7, & x=2.
    \end{cases}
  \]
  Calcule los límites laterales de $F_a$ cuando $x$ tiende a $2$ y determine
  todos los valores de $a$ para los cuales existe
  \[
    \lim_{x\to2}F_a(x).
  \]
  \textbf{Hint:} Factorice el numerador y considere los casos $x<2$ y $x>2$ por separado.
  \item \textbf{(3 puntos)} Sea $A$ una función definida para $x$ cercanos a $0$, excepto posiblemente en $0$, que satisface
  \[
    |A(x)|\leq3.
  \]
  Calcule mediante el teorema del Sanduche
  \[
    \lim_{x\to0}
    \frac{x^2}{|x|+x^2}\,A(x).
  \]
\textbf{Hint:} Encuentre una función $g$ tal que
  \[
    \left|\frac{x^2}{|x|+x^2}\,A(x)\right|\leq g(x)
  \]
  Note que $x^2=|x|^2$ para $x\neq 0$ y que $|xy|=|x||y|$ para todo $x,y\in\mathbb{R}$.
  \item \textbf{(4 puntos)} Considere las funciones
  \[
    f\colon\mathbb{R}\setminus\{2\}\to\mathbb{R}\setminus\{1\},
    \qquad
    f(x)=\frac{x+1}{x-2},
  \]
  y
  \[
    g\colon[0,\infty)\to[0,\infty),
    \qquad
    g(x)=\sqrt{x}.
  \]
  \begin{enumerate}
    \item \textbf{(2 puntos)} Analice la inyectividad y sobreyectividad de $f$.
    \item \textbf{(2 puntos)} Determine la función $g\circ f$.
  \end{enumerate}
\end{enumerate}

\ifmostrarSoluciones
\begin{solucion}[de la propuesta paralela]
\begin{enumerate}
  \item El numerador se factoriza como
  \[
    x^2+ax-2a-4=(x-2)(x+a+2).
  \]
  Si $x<2$, entonces $|x-2|=-(x-2)$ y
  \[
    F_a(x)=-(x+a+2).
  \]
  Si $x>2$, entonces $|x-2|=x-2$ y
  \[
    F_a(x)=x+a+2.
  \]
  Por tanto,
  \[
    \lim_{x\to2^-}F_a(x)=-(a+4),
    \qquad
    \lim_{x\to2^+}F_a(x)=a+4.
  \]
  El límite bilateral existe si y solo si
  \[
    -(a+4)=a+4,
  \]
  de donde $a=-4$. En ese caso,
  \[
    \lim_{x\to2}F_{-4}(x)=0.
  \]

  \item Para $x\neq0$, se tiene $x^2=|x|^2$ y $|x|>0$. Por tanto,
  \[
    \frac{x^2}{|x|+x^2}
    =\frac{|x|^2}{|x|(1+|x|)}
    =\frac{|x|}{1+|x|}.
  \]
  Como este factor es no negativo y $|A(x)|\leq3$, obtenemos
  \[
    \left|
      \frac{x^2}{|x|+x^2}A(x)
    \right|
    \leq
    3\frac{x^2}{|x|+x^2}
    =3\frac{|x|}{1+|x|}.
  \]
  En consecuencia,
  \[
    -\frac{3|x|}{1+|x|}
    \leq
    \frac{x^2}{|x|+x^2}A(x)
    \leq
    \frac{3|x|}{1+|x|}.
  \]
  Como $|x|$ tiende a $0$ y $1+|x|$ tiende a $1$, ambas cotas tienden a $0$.
  El teorema del Sanduche permite concluir que
  \[
    \lim_{x\to0}
    \frac{x^2}{|x|+x^2}A(x)=0.
  \]

  \item Para $x\neq2$, se tiene
  \[
    f(x)=1+\frac{3}{x-2}\neq1,
  \]
  por lo que $f$ toma valores en su codominio. Si $f(x_1)=f(x_2)$,
  al multiplicar por los denominadores no nulos obtenemos
  \[
    (x_1+1)(x_2-2)=(x_2+1)(x_1-2)
    \implies -3x_1=-3x_2
    \implies x_1=x_2.
  \]
  Así, $f$ es inyectiva. Para cada $y\neq1$, la ecuación $f(x)=y$ equivale a
  \[
    x=\frac{2y+1}{y-1}.
  \]
  Este valor pertenece al dominio, pues $x-2=3/(y-1)\neq0$, y al sustituirlo
  se verifica que $f(x)=y$. Por tanto, $f$ es sobreyectiva y biyectiva.

  La composición está definida cuando $x\neq2$ y $f(x)\geq0$, es decir,
  \[
    \frac{x+1}{x-2}\geq0.
  \]
  Los puntos críticos son $-1$ y $2$. El análisis de signos da
  \[
    \begin{array}{c|ccccc}
      x & (-\infty,-1) & -1 & (-1,2) & 2 & (2,\infty)\\ \hline
      (x+1)/(x-2) & + & 0 & - & \text{no definido} & +
    \end{array}
  \]
  Se incluye $x=-1$, pues la raíz de cero existe, y se excluye $x=2$,
  pues anula el denominador. En consecuencia,
  \[
    D_{g\circ f}=(-\infty,-1]\cup(2,\infty).
  \]
  Con este dominio natural, la función solicitada es
  \[
    g\circ f\colon(-\infty,-1]\cup(2,\infty)\to[0,\infty),
    \qquad
    (g\circ f)(x)=\sqrt{\frac{x+1}{x-2}}.
  \]
\end{enumerate}
\end{solucion}
\fi

\end{document}
